\documentclass[9pt,twocolumn,twoside]{pnas-new}
% Use the lineno option to display guide line numbers if required.

\articletype{CLASSIFICATION}%%%% Article topic/classification

\templatetype{pnasresearcharticle} % Choose template
% {pnasresearcharticle} = Template for a two-column research article
% {pnasmathematics} %= Template for a one-column mathematics article
% {pnasinvited} %= Template for a PNAS invited submission


\begin{document}


\title{Thermal niche warming is more consistent than range shifts in marine species under climate change}

% Use letters for affiliations, numbers to show equal authorship (if applicable) and to indicate the corresponding author
\author[a,1]{Federico Maioli}
\author[a]{P.\,Daniël van Denderen} 
\author[b]{Max Lindmark}
\author[a]{Marcel Montanyès}
\author[c]{Eric\,J. Ward}
\author[d,e]{Sean\,C. Anderson}
\author[a]{Martin Lindegren}




\affil[a]{National Institute of Aquatic Resources, Technical University of Denmark (DTU), Henrik Dams Allé 202, Kongens Lyngby, 2800, Denmark}  
\affil[b]{Department of Aquatic Resources, Swedish University of Agricultural Sciences (SLU), Turistgatan 5, Lysekil, 45330, Sweden} 
\affil[c]{Conservation Biology Division, NOAA Fisheries, Northwest Fisheries Science Center, 2725 Montlake Blvd E, Seattle, WA 98112, USA}  
\affil[d]{Pacific Biological Station, Fisheries and Oceans Canada, 3190 Hammond Bay Rd, Nanaimo, BC V9T 6N7, Canada}
\affil[e]{School of Resource and Environmental Management, Simon Fraser University, Burnaby, BC V5A 1S6, Canada}


% Please give the surname of the lead author for the running footer
\leadauthor{Maioli}

% Please add a significance statement to explain the relevance of your work
\significancestatement{Authors must submit a significance statement between 50 and 120 words in length about the significance of their research paper written at a level understandable to an undergraduate educated scientist outside their field of speciality. The primary goal of the significance statement is to explain the relevance of the work in broad context to a broad readership. If submitting a Direct Submission, please add a significance statement to explain the relevance of your work. Brief Reports do not publish with a significance statement, and should be omitted for this article type.}

% Please include corresponding author, author contribution and author declaration information
\authorcontributions{Please provide details of author contributions here.}
\authordeclaration{Please declare any competing interests here.}
\correspondingauthor{\textsuperscript{1}To whom correspondence should be addressed. E-mail: fedma@aqua.dtu.dk}

% At least three keywords are required at submission. Please provide three to five keywords, separated by the pipe symbol.
\keywords{Marine climate change $|$ Species redistribution $|$ Multidimensional range shifts $|$ Ocean warming}

\begin{abstract}
Marine species are widely expected to shift poleward or into deeper waters in response to rising ocean temperatures. However, our knowledge is primarily based on studies examining range shifts along single dimensions at a time (e.g., latitude or depth). Failing to address how movements along multiple dimensions interact, including associated changes in thermal exposure, may result in misleading conclusions and predictions of species distribution and community composition under global warming. To address this knowledge gap, we here develop and apply a multidimensional framework that jointly evaluates climate-driven redistribution of marine fish across latitude, longitude, depth and realized thermal niches, based on long-term scientific bottom-trawl surveys throughout the North Atlantic and Northeast Pacific. Our results show that net redistributions are generally small and highly region-specific, while the realized thermal niches of species have warmed substantially over the past three decades. These findings demonstrate that spatial redistribution is generally failing to keep pace with rising temperatures and challenge the prevailing assumption that marine species will move to escape warming. This has direct implications for biodiversity indicators that rely on distributional shifts as evidence of climate impacts, as well as climate-informed management and conservation of marine ecosystems, fisheries, and biodiversity at large.
\end{abstract}

\dates{This manuscript was compiled on \today}
\doi{\url{www.pnas.org/cgi/doi/10.1073/pnas.XXXXXXXXXX}}


\maketitle
\thispagestyle{firststyle}
\ifthenelse{\boolean{shortarticle}}{\ifthenelse{\boolean{singlecolumn}}{\abscontentformatted}{\abscontent}}{}

\Firstpage

%The \Firstpage command is used to format the first page text column size. The same size will be maintained for subsequent paragraph until the \Endparasplit or \Parasplit command is encountered.

The world's oceans are warming rapidly, driving a major reorganization of marine biodiversity with consequences for ecosystems and human use \citep{ipcc_ocean_2022, pinsky_preparing_2018, palacios-abrantes_climate_2025, pecl_biodiversity_2017}. Marine species are widely expected to respond to ocean warming by tracking their thermal niches---the temperature ranges suitable for growth, survival and reproduction \citep{grinnell_niche-relationships_1917}. Such tracking often involves shifts in geographic distribution \citep{tingley_birds_2009}. Consistent with this expectation, studies across marine taxa and regions document widespread poleward and deeper movements, making these range shifts among the most frequently reported biological responses to climate change \citep{parmesan_globally_2003, parmesan_ecological_2006, poloczanska_global_2013, poloczanska_responses_2016}.

Despite widespread reports of poleward and deeper shifts, some populations remain stationary even under sustained warming, while others shift in unexpected directions \citep{rubenstein_climate_2023, fuchs_wrong-way_2020, lawlor_mechanisms_2024}. These departures likely reflect the complex spatial structure of ocean warming, including heterogeneous climate velocities and physical constraints imposed by coastlines, bathymetry, and circulation patterns \citep{burrows_pace_2011, pinsky_marine_2013}. They may also arise from biological variability among taxa, including differences in thermal physiology, habitat affinity, recruitment capacity, and food-web limitations \citep{sunday_species_2015, roberts_substrate-dependent_2020, britten_changing_2016, zhang_food-web_2017}. What remains unclear is whether, once aggregated across species and regions, these diverse responses give rise to coherent net redistribution signal---such as poleward or deeper shifts---or whether no consistent pattern emerges under ocean warming.

\Endparasplit

Determining whether species responses combine into coherent redistribution patterns requires evaluating movement across multiple dimensions simultaneously. However, most previous studies examined redistribution along a single gradient at a time---most commonly latitude, or, less often, depth \citep{lenoir_climate-related_2015}. While these approaches have revealed important climate-related patterns and responses, they provide limited insight into how movements along different dimensions interact, and what this implies for net redistribution across regions. In marine systems, species may shift horizontally, move vertically into cooler waters, remain in place while tolerating greater thermal exposure, or combine these strategies \citep{nye_changing_2009}. Hence, without integrating these responses---including associated changes in realized niches---we lack a complete picture of climate-driven redistribution under ocean warming \citep{lawlor_mechanisms_2024, fredston_reimagining_2025}.

To address this gap, we develop and apply a multidimensional framework that evaluates climate-driven redistribution across three complementary dimensions: horizontal range shifts (changes in range centroids), vertical redistributions (realized depth niches), and changes in realized thermal niches. Unlike previous efforts that have modeled these attributes separately, our novel approach allows us to examine correlations across these redistribution pathways. This framework allows us to assess not only the direction and magnitude of spatial movements, but also whether such movements effectively track shifting thermal environments. If populations readily track changing isotherms, their realized thermal niches should remain stable through time. Hence, systematic warming of realized niches indicate incomplete tracking and increasing thermal exposure. We apply this framework to decades of standardized scientific bottom-trawl surveys encompassing more than 200 well-sampled demersal (bottom-living) fish populations across 11 regions of the North Atlantic and Northeast Pacific. These regions span broad latitudinal and climatic gradients, include some of the fastest-warming continental shelf seas \citep{pershing_slow_2015, song_arctic_2023}, and contain some of the most comprehensive long-term records of marine biodiversity change (Fig.~\ref{fig:map}; \citep{maureaud_are_2021, maureaud_fishglob_data_2024}). Using spatiotemporal and Bayesian mixed-effects models, we jointly estimate long-term trends in species’ horizontal distributions, depth niches, and realized thermal niches. By examining these responses both within and across regions, we assess whether species-level responses combine into coherent redistribution patterns across species and scales and evaluate the extent to which such movements reduce thermal exposure under ocean warming.



\section*{Guide to using this template on Overleaf}

Please note that whilst this template provides a preview of the typeset manuscript for submission, to help in this preparation, it will not necessarily be the final publication layout. For more detailed information please see the  \href{https://www.pnas.org/author-center/submitting-your-manuscript}{PNAS Author Center}.


If you have a question while using this template on Overleaf, please use the Menu on the top bar and access the Documentation for \href{https://www.overleaf.com/help}{help and tutorials}. You can also \href{https://www.overleaf.com/contact}{contact the Overleaf support team} at any time with specific questions.

\subsection*{Author Affiliations}

For each author, include institutional unit (e.g., division, department, or section), institution, city, state with ZIP code (for US institutions) or country with postal code (for non-US institutions). Use lower case letters to match authors with institutions, as shown in the example. PNAS requires the corresponding author to provide an ORCID identifier at submission and strongly encourages all authors to use an ORCID ID. Do not include ORCIDs in the manuscript file; individual authors must link their ORCID account to their PNAS profile at \href{http://www.pnascentral.org/}{www.pnascentral.org}. For proper authentication, authors must provide their ORCID at submission and are not permitted to add ORCIDs on proofs.

\subsection*{Submitting Manuscripts}

All authors must submit their articles at \href{http://www.pnascentral.org/cgi-bin/main.plex}{PNAScentral}. If you are using Overleaf to write your article, you can use the ``Submit to PNAS'' option in the top bar of the editor window.

\subsection*{Format}

Many authors find it useful to organize their manuscripts with the following order of sections: introduction, results, discussion, materials and methods, acknowledgments, and references. Other orders and headings are permitted.

\subsection*{Manuscript Length}

A standard 6-page article is approximately 4,000 words, 50 references, and 4 medium-size graphical elements (i.e., figures and tables). The preferred length of Direct Submission articles remains at 6 pages, but PNAS will allow articles up to a maximum of 12 pages. Brief Reports are limited to 3 pages, which is approximately 1,600 words (including the manuscript text, title page, abstract, and figure legends), and 15 references.

\subsection*{References}

References should be cited in numerical order as they appear in text; this will be done automatically via bibtex, e.g. \cite{belkin2002using} and \cite{berard1994embedding,coifman2005geometric,phdthesis,masterthesis}. All references cited in the main text should be included in the main manuscript file.\Endparasplit

%The \Endparasplit command is used to end the formatting of the text column size that was set by the \Firstpage command. This will restore the default column size for the subsequent text.
%
%The \Parasplit command should be used if the page ends in the middle of a paragraph, allowing the column formatting to continue seamlessly on the next page.


\subsection*{Language-Editing Services}
Prior to submission, authors who believe their manuscripts would benefit from professional editing are encouraged to use a language-editing service (see list at https://www.pnas.org/author-center/language-editing). PNAS does not take responsibility for or endorse these services, and their use has no bearing on acceptance of a manuscript for publication.


\begin{figure}%[tbhp]
\centering
\includegraphics[width=.8\linewidth]{frog.pdf}
\caption{Placeholder image of a frog with a long example legend to show justification setting.}
\label{fig:frog1}
\end{figure}


\begin{figure*}[t!]
\centering
\includegraphics[scale=0.6]{frog}
\caption{Placeholder image of a frog with a long example legend to show justification setting.}
\label{fig:frog2}
\end{figure*}


\begin{table}[t!]
\centering
\caption{Comparison of the fitted potential energy surfaces and ab initio benchmark electronic energy calculations}
\begin{tabular}{lrrr}
Species & CBS & CV & G3 \\
\midrule
1. Acetaldehyde & 0.0 & 0.0 & 0.0 \\
2. Vinyl alcohol & 9.1 & 9.6 & 13.5 \\
3. Hydroxyethylidene & 50.8 & 51.2 & 54.0\\
\bottomrule
\end{tabular}

\addtabletext{nomenclature for the TSs refers to the numbered species in the table.}
\end{table}


\begin{table*}[t!]
\centering
\caption{Impact on Emission Behaviors by Socioeconomic Status}
\begin{tabular*}{\textwidth}{@{\extracolsep{\fill}}lccccc@{}}
& (1) & (2) & (3) & (4) & (5) \\
\midrule
 & \multicolumn{2}{c}{City-level} & Firm-level & \multicolumn{2}{c}{City-level} \\
Dep. Var.: & \multicolumn{2}{c}{COD Emission (1,000 tons)} & COD Emission (ton) & Firm Entry & Firm Exit \\
\multicolumn{6}{c}{Panel A: Population Share without College Education} \\
Share of Below College $\times$ Post$_{05}$ & 0.165*** & 0.292** & 0.358*** & 0.623** & 0.280 \\
 & (0.026) & (0.119) & (0.106) & (0.266) & (0.487) \\
Share of Below College & -0.208 & & & & \\
 & (0.141) & & & & \\
Post$_{05}$ & -16.426*** & & & & \\
 & (2.873) & & & & \\
\multicolumn{6}{c}{Panel B: Population Share without High School Education} \\
Share of Below HS $\times$ Post$_{05}$ & 0.099*** & 0.218** & 0.232** & 0.453** & 0.089 \\
 & (0.005) & (0.090) & (0.079) & (0.195) & (0.333) \\
Share of Below HS & -0.213* & & & & \\
 & (0.100) & & & & \\
Post$_{05}$ & -9.465*** & & & & \\
 & (0.564) & & & & \\
\bottomrule
\end{tabular*}

\addtabletext{*** P $<$ 0.01, ** P $<$ 0.05, * P $<$ 0.1}
\end{table*}

\subsection*{Digital Figures}

EPS and high-resolution PDF are preferred formats for figures that will be used in the main manuscript. Authors may submit PRC or U3D files for 3D images; these must be accompanied by 2D representations in TIFF, EPS, or high-resolution PDF format. Color images must be in RGB (red, green, blue) mode. Include the font files for any text.

Images must be provided at final size, preferably 1 column width (8.7cm). Figures wider than 1 column should be sized to 11.4cm or 17.8cm wide. Numbers, letters, and symbols should be no smaller than 6 points (2mm) and no larger than 12 points (6mm) after reduction and must be consistent.

Figures and tables should be labelled and referenced in the standard way using the \verb|\label{}| and \verb|\ref{}| commands.

Figure \ref{fig:frog1} shows an example of how to insert a column-wide figure. To insert a figure wider than one column, please use the \verb|\begin{figure*}...\end{figure*}| environment. Figures wider than one column should be sized to 11.4 cm or 17.8 cm wide. Use \verb|\begin{SCfigure*}...\end{SCfigure*}| for a wide figure with side legends.

\begin{SCfigure*}[\sidecaptionrelwidth][t!]
\centering
\includegraphics[width=11.4cm,height=11.4cm]{frog.pdf}
\caption{This legend would be placed at the side of the figure, rather than below it.}\label{fig:side}
\end{SCfigure*}


\subsection*{Tables}
Tables should be included in the main manuscript file and should not be uploaded separately.


\subsection*{Footnotes}

Footnotes are ordered by symbols such as *, \textdagger, \textdaggerdbl, etc. The article should contain only 13 footnotes\footnote{Footnote Example 1}\footnote{Footnote Example 2}.


\subsection*{Single column equations}

Authors may use 1- or 2-column equations in their article, according to their preference.

\begin{align*}
(x+y)^3&=(x+y)(x+y)^2\\
       &=(x+y)(x^2+2xy+y^2) \numberthis \label{eqn:example1} \\
       &=x^3+3x^2y+3xy^3+x^3.
\end{align*}

To allow an equation to span both columns, use the \verb|\begin{figure*}...\end{figure*}| environment mentioned above for figures.

Note that the use of the \verb|widetext| environment for equations is not recommended, and should not be used.

\begin{figure*}[bt!]
\begin{align*}
(x+y)^3&=(x+y)(x+y)^2\\
       &=(x+y)(x^2+2xy+y^2) \numberthis \label{eqn:example2} \\
       &=x^3+3x^2y+3xy^3+x^3.
\end{align*}
\end{figure*}



\subsection*{Supporting Information Appendix (SI)}

Authors should submit SI as a single separate SI Appendix PDF file, combining all text, figures, tables, movie legends, and SI references. SI will be published as provided by the authors; it will not be edited or composed. Additional details can be found in the \href{https://www.pnas.org/author-center/submitting-your-manuscript#supporting-information}{PNAS Author Center}. The PNAS Overleaf SI template can be found \href{https://www.overleaf.com/latex/templates/pnas-template-for-supplementary-information/wqfsfqwyjtsd}{here}. Refer to the SI Appendix in the manuscript at an appropriate point in the text. Number supporting figures and tables starting with S1, S2, etc.

Authors who place detailed materials and methods in an SI Appendix must provide sufficient detail in the main text methods to enable a reader to follow the logic of the procedures and results and also must reference the SI methods. If a paper is fundamentally a study of a new method or technique, then the methods must be described completely in the main text.

Supporting information (SI) for PNAS Brief Reports is limited to extended methods, essential supporting datasets, software, and video/audio files; supporting text, tables or figures are not permitted.

\subsubsection*{SI Datasets}

Supply .xlsx, .csv, .txt, .rtf, GZ or .pdf files. This file type will be published in raw format and will not be edited or composed.


\paragraph*{SI Movies}

Supply Audio Video Interleave (avi), Quicktime (mov), Windows Media (wmv), animated GIF (gif), or MPEG-4 Part 14 (mp4) files. Movie legends should be included in the SI Appendix file. All movies should be submitted at the desired reproduction size and length. Movies should be no more than 10MB in size.



\matmethods{Please describe your materials and methods here. This can be more than one paragraph, and may contain subsections and equations as required. If your research involved human or animal participants, please identify the institutional review board and/or licensing committee that approved the experiments. Please also include a brief description of your informed consent procedure if your experiments involved human participants.


\subsection*{Subsection for Method}
Example text for subsection.
}

\showmatmethods{} % Display the Materials and Methods section

\dataavail{To allow others to replicate and build on work published in PNAS, authors must make materials, data, and associated protocols, including code and scripts, available to readers in a public repository upon publication. Restrictions on full or partial access to these materials and requests for legal, ethical, and logistical (e.g., size) exceptions must be noted at submission. If requested, these materials must be made available to editors and reviewers during submission for the purpose of evaluating the manuscript. A statement detailing sharing plans will be included in the published article as provided within the submission form. Research datasets, whether original or previously published, must be cited in the references as a condition for publication. Please refer to our full \href{https://www.pnas.org/author-center/editorial-and-journal-policies\#materials-and-data-availability}{policy}.}

\acknow{We thank all those who contributed to the collection of the survey data and the FISHGLOB consortium for providing public access to the dataset and for discussions held in Oldenburg. We thank Chris Rooper for comments that improved the manuscript. F.M. also thanks Christopher Griffiths for early discussions that helped shape this work.}

\showacknow{} % Display the acknowledgments section




\bibsplit[3]
%Use \bibsplit to split the references from the body of the text. Value "[3]" represents the number of reference in the left column (Note: Please avoid single column figures & tables on this page.)

% Bibliography
\bibliography{references}

\end{document}
